%Problem
\section{Problem Formulation}

This section presents a mathematical model of the Resource Constrained Capacitated Vehicle Routing Problem with Time Windows (RCVRPTW). The model incorporates multiple operational constraints including vehicle capacities, time windows, service times, and working time limits.

\subsection{Problem Definition}

We assume the existence of $N$ service locations in set $\mathcal{N} = \{1, 2, \dots, N\}$. This set is extended by location 0, representing the depot, creating a set $\mathcal{D} = \{0\} \cup \mathcal{N}$. The depot is the start and end point for each service vehicle. Each location has defined coordinates, based on which the travel cost between locations is calculated. The cost is represented by the distance matrix $d_{i,j}$ denoting the travel cost from location $i$ to $j$.

A fleet of vehicles $\mathcal{V}$ is available, where each vehicle $v \in \mathcal{V}$ has a capacity constraint $Q_v$. Each customer location $i \in \mathcal{N}$ has a demand $q_i$, a service time $s_i$, and a time window $[a_i, b_i]$ during which service should ideally occur. The maximal capacity is set to 1000~kg, which represents an average load for small and medium delivery vans. The time constraint is set to 480 minutes, consistent with an 8-hour working day.

The notation used throughout this section is summarized in Table~\ref{tab:notation}.

\begin{table}[ht]
	\renewcommand{\arraystretch}{1.3}
	\caption{Notation summary}
	\label{tab:notation}
	\centering
	\begin{tabular}{l|l}
		\hline\hline
		Symbol & Meaning  \\
		\hline \hline
		$\mathcal{V}$ & Set of vehicles \\ 
        \hline
        $\mathcal{D}$ & Set of all locations (including depot) \\
        \hline
        $\mathcal{N}$ & Set of customer locations ($\mathcal{D}\setminus\{0\}$) \\
        \hline
        $d_{i,j}$ & Travel cost from location $i$ to $j$ \\
        \hline
        $s_i$ & Service time at location $i$ \\
		\hline
        $[a_i,b_i]$ & Time window for location $i$\\
        \hline
        $q_i$ & Demand at customer location $i$\\
        \hline
        $Q_v$ & Capacity of vehicle $v$ \\
        \hline
        $x_{v,i,j}$ & Binary variable: 1 if vehicle $v$ uses arc $i\rightarrow j$ \\
        \hline
        $z_{v,i}$ & Binary variable: 1 if vehicle $v$ serves location $i$\\
        \hline
        $A_{v,i}$ & Arrival time of vehicle $v$ at location $i$\\
        \hline
        $S_{v,i}$ & Service start time of vehicle $v$ at location $i$\\
        \hline
        $w_{v,i}$ & Waiting time of vehicle $v$ at location $i$ ($S_{v,i}-A_{v,i}$)\\
        \hline
        $p_{v,i}$ & Early arrival penalty for vehicle $v$ at $i$ \\ 
        \hline
        $q_{v,i}$ & Late arrival penalty for vehicle $v$ at $i$ \\ 
        \hline
        $y_{v,i}$ & Load of vehicle $v$ after servicing location $i$\\
        \hline
        $\alpha$ & Penalty coefficient for time window violations\\
        \hline
        $M$ & Large constant (big-M method)\\
        \hline \hline
	\end{tabular}
\end{table}

\subsection{Mathematical Model}

The mathematical model consists of an objective function that minimizes the total operational cost, subject to a set of constraints that ensure feasibility. The objective function comprises three main components: travel costs, waiting time costs, and penalty costs for time window violations.

\subsubsection{Objective Function}

The complete objective function minimizes the total cost:
\begin{equation}\label{eq:objective}
     \min \left( \sum_{v \in \mathcal{V}} \sum_{i \in \mathcal{D}} \sum_{j \in \mathcal{D}}  x_{v,i,j} d_{i,j} + \sum_{v \in \mathcal{V}}\sum_{i \in \mathcal{D}} w_{v,i} + \sum_{v \in \mathcal{V}}\sum_{i \in \mathcal{D}} \alpha (p_{v,i} + q_{v,i}) \right)
\end{equation}

where the first term represents the total travel distance, the second term accounts for waiting time costs (vehicle and driver idle time), and the third term penalizes time window violations with penalty coefficient $\alpha$.

\subsubsection{Constraints}

The objective function \eqref{eq:objective} is minimized subject to the following constraints:

\textbf{Route structure constraints:}

\begin{equation}{\label{eq:one_visit}}
     \sum_{v \in \mathcal{V}}\sum_{j \in \mathcal{D}\setminus \{i\}} x_{v,i,j}=1, \quad \forall i \in \mathcal{N}
\end{equation}

\begin{equation}{\label{eq:flow_balance}}
    \sum_{j\in\mathcal{D}\setminus\{i\}}x_{v,i,j}-\sum_{j\in\mathcal{D}\setminus\{i\}}x_{v,j,i}=0, \quad \forall v \in \mathcal{V}, i \in \mathcal{N}
\end{equation}

\begin{equation}{\label{eq:start_end_in_depot}}
    \sum_{j\in\mathcal{N}}x_{v,0,j}=z_v, \quad \sum_{i\in\mathcal{N}}x_{v,i,0}=z_v, \quad \forall v \in \mathcal{V}
\end{equation}

\begin{equation}{\label{eq:edgevisitcorrelation}}
    z_{v,i} = \sum_{j\in \mathcal{D}\setminus\{i\}}x_{v,j,i} = \sum_{j\in \mathcal{D}\setminus\{i\}}x_{v,i,j}, \quad \forall v \in \mathcal{V}, i \in \mathcal{N}
\end{equation}

Constraint \eqref{eq:one_visit} ensures that each customer location is visited exactly once by one vehicle. Constraint \eqref{eq:flow_balance} enforces flow conservation: if a vehicle arrives at a location, it must also depart from it. Constraint \eqref{eq:start_end_in_depot} requires that each vehicle starts and ends at the depot. Constraint \eqref{eq:edgevisitcorrelation} links the visit indicator variable with the routing variables.

\textbf{Time window constraints:}

\begin{equation}{\label{eq:time_cohestion}}
 A_{v,j} \;\ge\; S_{v,i} + s_i + d_{i,j} \;-\; M(1 - x_{v,i,j}), \quad \forall v \in \mathcal{V}, \forall i,j \in \mathcal{D}, i\neq j
\end{equation}

\begin{equation}\label{eq:start_soft}
S_{v,i} \;\ge\; A_{v,i}, \qquad a_i \;\le\; S_{v,i} + p_{v,i}, \quad p_{v,i}\ge 0, \quad \forall v \in \mathcal{V}, i \in \mathcal{N}
\end{equation}

\begin{equation}\label{eq:penalty}
S_{v,i} \;\le\; b_i + q_{v,i}, \quad q_{v,i}\ge 0, \quad \forall v \in \mathcal{V}, i \in \mathcal{N}
\end{equation}

\begin{equation}\label{eq:waiting_time}
w_{v,i} \;=\; S_{v,i} - A_{v,i}, \quad w_{v,i}\ge 0, \quad \forall v \in \mathcal{V}, i \in \mathcal{N}
\end{equation}

\begin{equation}{\label{eq:notimesforvehiclenotvisit}}
    A_{v,i}\leq M z_{v,i}, \quad S_{v,i}\leq M z_{v,i}, \quad \forall v \in \mathcal{V}, i \in \mathcal{N}
\end{equation}

Constraint \eqref{eq:time_cohestion} ensures temporal consistency: arrival at the next location cannot occur before completing service at the current location plus travel time. Constraints \eqref{eq:start_soft} and \eqref{eq:penalty} implement soft time windows by introducing penalty variables $p_{v,i}$ (early arrival) and $q_{v,i}$ (late arrival). Constraint \eqref{eq:waiting_time} defines waiting time as the difference between service start and arrival times. Constraint \eqref{eq:notimesforvehiclenotvisit} deactivates time variables for locations not visited by vehicle $v$.

\textbf{Capacity constraints:}

\begin{equation}\label{eq:capacity_flow}
y_{v,j} \;\ge\; y_{v,i} + q_j - Q_v(1 - x_{v,i,j}), \quad \forall v \in \mathcal{V}, i,j \in \mathcal{D}, i \neq j
\end{equation}

\begin{equation}\label{eq:capacity_limits}
0 \;\le\; y_{v,i} \;\le\; Q_v z_{v,i}, \quad \forall v \in \mathcal{V}, i \in \mathcal{N}, \qquad y_{v,0} = 0, \forall v \in \mathcal{V}
\end{equation}

Constraint \eqref{eq:capacity_flow} tracks cumulative vehicle load: when traveling from $i$ to $j$, the load increases by the demand at $j$. Constraint \eqref{eq:capacity_limits} ensures that the vehicle load remains non-negative and does not exceed vehicle capacity.

\textbf{Variable domains:}

\begin{equation}
     x_{v,i,j}, z_{v,i}\in\{0,1\}, \quad \forall v \in \mathcal{V}, i,j \in \mathcal{D}
\end{equation}

\begin{equation}
     p_{v,i}, q_{v,i}, w_{v,i}, A_{v,i}, S_{v,i}, y_{v,i} \geq 0, \quad \forall v \in \mathcal{V}, i \in \mathcal{D}
\end{equation}

This formulation provides a complete Mixed Integer Linear Programming (MILP) model for the RCVRPTW that can be solved using commercial solvers or tackled with metaheuristic approaches for large-scale instances, as described in the subsequent sections.
